\section{Perception and Color}

Visualization works because it offloads cognition onto the human \term{visual
system}, a high-bandwidth, partly pre-cognitive channel. Knowing how the eye
and early vision behave tells us \emph{which} visual channels are fast,
accurate and orderable, and which ones (notably hue) are not. This section
grounds the design rules used elsewhere (expressiveness, effectiveness, channel
ranking) in perception, and then covers color, the single most over-used and
most-often-misused channel in the exam's ``Visualization Critique'' questions.

\begin{keybox}[Why perception matters for vis]
A good encoding lets the visual system answer the task \emph{pre-attentively}
(in a glance, $\sim$200\,ms) instead of forcing slow serial reading.
Bad encodings (rainbow color, 3D bars, hue for ordered data) fight the visual
system and force slow, error-prone interpretation.
\end{keybox}

\subsection{The Human Visual System}

\subsubsection{Eye, retina, rods and cones}
Light from the scene is projected by the cornea/lens onto the \term{retina} at
the back of each eye. The retina is covered with photoreceptors:
\begin{itemize}
  \item \term{Cones} -- color-sensitive (three types, tuned to long / medium /
    short wavelengths), concentrated in the centre of the retina; responsible
    for detailed, color daytime vision.
  \item \term{Rods} -- highly light-sensitive but \emph{not} color-sensitive;
    dominate peripheral and low-light vision.
\end{itemize}

\subsubsection{Fovea and acuity}
The \term{fovea} is the small central region of the retina where image detail
and color rendition are best (densest cone packing). Crucially, the fovea
covers a viewing angle of only about \textbf{$2^\circ$} -- roughly a thumbnail
held at arm's length. Moving outward through the \term{parafovea} to the
\term{periphery}, \term{visual acuity} (the ability to resolve detail)
\emph{falls off dramatically} [Spence, Sec.~4.1].

\subsubsection{Fixations and saccades}
Because only a tiny region is seen in detail, the eye constantly jumps around
the scene. Each \term{fixation} (gaze held on one point) lasts about
\textbf{200\,ms}; fixations are connected by very fast ($\sim$20\,ms)
\term{saccades}. Viewing a whole image involves $\sim$100 fixations.
\term{Overt attention} is paid to a few fixated points; \term{covert
attention} to parafoveal/peripheral items helps decide where to look next.

\subsubsection{Visual cortex (it's complicated)}
Retinal signals pass through the LGN to the primary visual cortex \term{V1}
(the main entry point; neurons highly selective for gradients, i.e.\ ``edge
detection''), then V2/V4 with increasingly complex response, up to the inferior
temporal cortex (object recognition, memory). Exam-relevant takeaway:
\emph{luminance contrast drives edge detection}, which is why luminance is a
good channel for fine spatial detail and hue is not.

\sketchnote{Retina disk with a small bright dot labelled \emph{fovea}
($2^\circ$), a ring \emph{parafovea}, and an outer ring \emph{periphery}; an
arrow ``acuity'' pointing outward and shrinking. Beside it: a zig-zag gaze path
(fixations $\approx$200\,ms, saccades $\approx$20\,ms).}

\subsection{Preattentive Processing and Popout}

A \term{limited} set of visual channels is detected \emph{rapidly and in
parallel}, within $\sim$200--250\,ms, before focused attention engages.

\begin{keybox}[Operational definition (exam-quotable)]
If finding a target takes a \textbf{fixed amount of time regardless of the
number of distractors}, the channel is processed \term{preattentively}.
The target then ``pops out''.
\end{keybox}

A single red dot among blue dots pops out instantly whether there are 12 or 50
distractors. Channels that support popout include color (hue), orientation
(tilt), size, curvature, motion, and others (Healey's catalogue).

\begin{pitfall}[Conjunction search -- the key limitation]
Popout works only when the target differs in \emph{one} unique visual property.
If the target shares each of its properties with some distractors and is
defined only by a \term{combination} of features (e.g.\ ``the red \emph{circle}''
among red squares and blue circles), there is no unique channel to search for.
This is a \term{conjunction search}: \textbf{not possible pre-attentively}, so
a slow \emph{serial} search is required and search time grows with item count.
\end{pitfall}

\textbf{Important nuance:} a conjunction of \emph{spatial position} with one
other channel (color or shape) \emph{can} still be searched preattentively --
e.g.\ in a scatterplot the conjunction of space~$\times$~color, or
space~$\times$~shape, pops out. The hard case is conjunctions of two
non-spatial channels.

\begin{exambox}[How this is asked]
``Why is it slow to find X in this chart?'' $\Rightarrow$ identify whether the
target has a unique preattentive feature or requires conjunction search.
Design fix: encode the target's category with a \emph{single} popout channel
(e.g.\ make the outliers a distinct hue) so they pop out preattentively.
\end{exambox}

\subsection{Gestalt Laws of Grouping}

\term{Gestalt laws} describe how we perceive \term{patterns / groupings} in
visual displays, and translate directly into design principles for information
displays.

\begin{longtable}{@{}p{2.6cm}p{7.2cm}p{4.4cm}@{}}
\toprule
\textbf{Law} & \textbf{We group items that\ldots} & \textbf{Vis example} \\
\midrule
\endhead
Proximity     & are close together in space            & clusters in a scatterplot read as groups \\
Similarity    & look alike (same hue/shape/size)        & same-color points = same category \\
Connectedness & are joined by lines                     & node--link edges, slope graphs \\
Continuity    & lie along a smooth continuous contour   & a line is followed through crossings \\
Closure       & form a closed/implied contour           & we ``complete'' a shape \\
Symmetry      & are symmetric about an axis             & mirrored layouts read as one object \\
Common fate   & move together                           & animated points moving together group \\
Common region & lie inside a shared enclosure           & a shaded hull / bounding box \\
\bottomrule
\end{longtable}

\begin{keybox}[Which is strongest?]
\term{Connectedness} (and shared enclosure / common region) can be a
\emph{more powerful} grouping cue than \term{proximity}: items joined by a
line or surrounded by a hull are grouped even when other cues (color or
distance) would group them differently. The slides state explicitly:
\emph{``connectedness can be a more powerful grouping principle than
proximity.''}
\end{keybox}

\sketchnote{Two scatter clusters. (a) proximity alone $\to$ two groups by
distance. (b) draw a connecting line / shaded blob that links one point from
the left cluster to the right cluster $\to$ it now reads as belonging to the
enclosed group, overriding proximity.}

\subsection{Weber's Law and the Just Noticeable Difference}

\subsubsection{JND}
The \term{Just Noticeable Difference (JND)} is the amount a stimulus must be
changed to be detectable by human observers \emph{at least half of the time}.
It is determined \emph{empirically}.

\subsubsection{Weber's law}
The JND is \emph{not} an absolute amount: it is proportional to the background
stimulus magnitude. To perceive a difference between a background intensity
$I$ and the background plus stimulation $\Delta I$, the difference must be
proportional to the background:
\[
  k = \frac{\Delta I}{I} \quad\Longleftrightarrow\quad \Delta I = k\,I,
  \qquad k = \text{const (Weber fraction)}.
\]
\textbf{Implication:} we judge \emph{relative}, not absolute, differences.
Adding 5 dots to a field of 10 is obvious; adding 5 to a field of 100 is not,
even though the absolute change is identical.

\begin{keybox}[Consequences for design]
\begin{itemize}
  \item Integrating Weber's law over stimulus intensity gives \term{Fechner's
    law}: perceived intensity is \emph{logarithmic},
    $p = k\ln(I/I_0)$, where $I_0$ is the absolute threshold.
  \item More generally, \term{Stevens' power law} models perceived sensation
    as $p = k\,I^{\,a}$ ($a$ a measured exponent per channel). Length
    $a\!\approx\!1$ (perceived $\approx$ veridical); area $a\!\approx\!0.7$ and
    brightness $a\!\approx\!0.5$ are \emph{compressed} (under-estimated), while
    saturation and electric shock are exaggerated. This is the perceptual basis
    for why \emph{position/length beats area/color} in the channel ranking.
  \item Practical rule: length differences are most accurately judged when bars
    share a \textbf{common baseline and are adjacent}; offset or stacked bars
    make comparison much harder.
\end{itemize}
\end{keybox}

\subsection{Color}

\subsubsection{Color is three channels}
Colloquial ``color'' is best understood as \textbf{three separate channels}:
\term{luminance} (light--dark), \term{hue} (red, green, blue, \ldots), and
\term{saturation} (vividness). The major colormap design choice is which of
these to use for what.

\begin{keybox}[Ordered vs.\ unordered channels (Munzner)]
\begin{itemize}
  \item \term{Luminance}: a \emph{magnitude} (ordered) channel -- people
    reliably order grays, always placing gray between black and white.
  \item \term{Saturation}: ordered, but \emph{less easily} (only $\sim$3--4
    discriminable steps for small marks).
  \item \term{Hue}: an \emph{identity} (categorical, \textbf{unordered})
    channel. Asked to order red/blue/green/yellow patches, people disagree.
    \textbf{Hue cannot be ordered} -- learned conventions (rainbow, traffic
    lights) are higher-level, not perceptual.
\end{itemize}
\end{keybox}

\subsubsection{Opponent color theory}
The retina's signals are immediately processed into \textbf{three opponent
color channels}: red--green, blue--yellow, and a black--white channel encoding
\term{luminance}. The luminance channel carries \emph{high-resolution} edge
information; the two chromatic channels are \emph{lower resolution}. This is
why luminance (not hue) should carry fine spatial detail.

Opponent theory also explains \term{color blindness} (more accurately
\emph{color vision deficiency}), affecting $\sim$8--10\% of men in its common
form: the \term{red--green} channel is reduced/absent while blue--yellow still
works $\Rightarrow$ \textbf{avoid red--green encodings}.

\subsubsection{Colormap types and when to use each}
The colormap taxonomy mirrors the data-type taxonomy.

\begin{longtable}{@{}p{2.9cm}p{2.4cm}p{8.7cm}@{}}
\toprule
\textbf{Colormap} & \textbf{Data} & \textbf{Use / rule} \\
\midrule
\endhead
\term{Categorical} (qualitative) & categorical &
  Distinct hues for unordered groups; \emph{segmented}. Hue = identity channel.
  Limited to \textbf{$\sim$6--12 discriminable bins} (12 max). Requires a
  \textbf{legend}. \\
\term{Sequential} & ordered (ord./quant.) &
  Progression from min to max along an ordered channel (luminance, possibly
  with one hue), e.g.\ light$\to$dark pink. \\
\term{Diverging} & ordered with meaningful midpoint &
  Two sequential ramps meeting at a \textbf{neutral midpoint} (zero / mean),
  e.g.\ blue--white--red ``warming stripes''. Use when a zero/reference point
  is meaningful (gain vs.\ loss). \\
\term{Bivariate} & two attributes &
  Encodes two attributes at once. Easy to read when the 2nd attribute is
  \emph{binary}; quickly unreadable when both have several levels. \\
\bottomrule
\end{longtable}

Colormaps can be \term{continuous} (good for quantitative / spatial fields) or
\term{segmented}/binned (good for categorical; emphasizes discrete bins).

\begin{keybox}[Rules for categorical color]
\begin{itemize}
  \item Use \textbf{$\le$ 6--12} colors; more than 12 hues are very hard to
    discriminate (Ware). Count the background and default object colors
    (black/white/gray) in your budget.
  \item Prefer \textbf{easily nameable} colors (only $\sim$8 are consistently
    named): start with the opponent-axis colors red/green/blue/yellow, then
    orange/brown/pink/purple/cyan.
  \item Match saturation to mark size: \emph{small} marks (points, lines)
    $\to$ highly saturated; \emph{large} areas $\to$ low saturation/pastel.
  \item If categories exceed the color budget, aggregate into groups or an
    ``other'' bin, or switch to an additional channel (shape, size).
  \item Categorical color always needs a \textbf{legend}.
\end{itemize}
\end{keybox}

\subsubsection{The rainbow / jet colormap problem}
Using the rainbow (``jet'') colormap for \emph{ordered} data is a classic
mistake. Munzner lists three serious perceptual problems:

\begin{pitfall}[Three problems with the rainbow colormap]
\begin{enumerate}
  \item \textbf{Hue is used to indicate order}, but hue is an identity channel
    with \emph{no implicit perceptual ordering} (expressiveness mismatch).
  \item \textbf{Not perceptually uniform}: equal steps in data are \emph{not}
    perceived as equal -- some ranges show several distinct colors, others
    look uniform (a wide band that just ``looks yellow'').
  \item \textbf{Fine detail is lost}: hue cannot convey fine-grained detail;
    luminance would, because luminance contrast drives edge detection.
\end{enumerate}
Add a fourth, practical one: rainbow is \textbf{not color-blind safe}
(red/green confusions).
\end{pitfall}

\subsubsection{Perceptually uniform maps (viridis)}
The fixes: (i) design \term{monotonically increasing luminance} colormaps so
that order is carried by the magnitude (luminance) channel while hue only marks
semantic regions; (ii) use a \textbf{perceptually uniform} colormap such as
\term{viridis} -- built in a perceptually uniform space (cf.\ CIE
$L^*a^*b^*$) so equal data steps look equally different, with monotonic
luminance, and it is color-blind friendly. (A perceptually linear rainbow is
possible but loses brightness/dynamic range, so it is rarely used.)

\begin{exambox}[Color critique checklist]
For any color encoding, ask:
(1) Is the data \emph{ordered}? If yes, do \emph{not} use hue alone -- use
luminance / a sequential or diverging map.
(2) Is there a meaningful midpoint? $\to$ diverging.
(3) Categorical with $\le$12 groups? $\to$ distinct nameable hues + legend.
(4) Is it color-blind safe (no red--green dependence; redundant encoding)?
(5) Is it perceptually uniform (avoid rainbow/jet; prefer viridis)?
\end{exambox}

\begin{pitfall}[color $\neq$ color]
``Color'' is not one thing. Saying ``encode this with color'' is too vague:
state \emph{which} channel (luminance/hue/saturation) and check it matches the
attribute type (ordered $\to$ luminance; categorical $\to$ hue).
\end{pitfall}
